

\SGARefTarget{sga3:iv:section:3}\section*{3. Universal Effective Equivalence Relations}
\addcontentsline{toc}{section}{3. Universal Effective Equivalence Relations}

\SGARefTarget{sga3:iv:subsection:3:1}\subsection*{3.1. Equivalence Relations: Definitions}
\addcontentsline{toc}{subsection}{\SGARefLink{sga3:iv:subsection:3:1}{3.1}. Equivalence Relations: Definitions}

\noindent\SGARefTarget{sga3:iv:definition:3:1:1}\textbf{Definition 3.1.1.}
\textit{A \(\mathscr C\)-equivalence relation in an object
\(X\in\operatorname{Ob}\mathscr C\) is a representable subfunctor \(R\) of
the functor \(X\times X\) such that, for every
\(S\in\operatorname{Ob}\mathscr C\), \(R(S)\) is the graph of an equivalence
relation on \(X(S)\).}

\SGARefLink{sga3:iv:definition:3:1:1}{This definition} applies in particular to the category
\(\widehat{\mathscr C}\). If \(X\) is regarded as an object of
\(\widehat{\mathscr C}\), then a \(\widehat{\mathscr C}\)-equivalence relation
in \(X\) is simply a subfunctor \(R\) of \(X\times X\) such that \(R(S)\) is
the graph of an equivalence relation on \(X(S)\), for every object \(S\) of
\(\mathscr C\).

If \(R\) is a \(\mathscr C\)-equivalence relation in \(X\), let \(p_i\)
denote the morphism from \(R\) to \(X\) induced by the projection
\(\operatorname{pr}_i:X\times X\to X\). Thus there is a diagram
\[
  p_1,p_2:R\rightrightarrows X.
\]

\medskip
\noindent\SGARefTarget{sga3:iv:definition:3:1:2}\textbf{Definition 3.1.2.}
\textit{A morphism \(u:X\to Z\) is said to be \emph{compatible with \(R\)}
if \(up_1=up_2\). An object of \(\mathscr C\) that is a cokernel of the pair
\((p_1,p_2)\) is also called a \emph{quotient object} of \(X\) by \(R\), and
is denoted \(X/R\).}

Thus one has an exact diagram
\[
  R\overset{p_1}{\underset{p_2}{\rightrightarrows}}
  X\xrightarrow{p}X/R,
\]
and \(X/R\) represents the covariant functor
\[
  \operatorname{Hom}_{\mathscr C}(X/R,Z)
  =\{\text{morphisms from \(X\) to \(Z\) compatible with \(R\)}\}.
\]
Since quotient objects are taken in \(\mathscr C\), the quotient \(X/R\) is
unique whenever it exists.

These definitions extend at once to a
\(\widehat{\mathscr C}\)-equivalence relation in \(X\). One should note,
however, that the identification of the objects of \(\mathscr C\) with their
images in \(\widehat{\mathscr C}\) does not commute with the formation of
quotients: the quotient \(X/R\) of \(X\) by \(R\) in \(\mathscr C\) is not
\emph{a priori} a quotient of \(X\) by \(R\) in
\(\widehat{\mathscr C}\). One must therefore avoid identifying
\(\mathscr C\) uncritically with its image in \(\widehat{\mathscr C}\) in
questions involving passage to a quotient.\footnote{Editor's note: here is an
illustration that also previews the sequel of this expos\'e. Let \(G\) be a
\(\mathscr C\)-group, let \(H\) be a \(\mathscr C\)-subgroup, and let \(R\)
be the equivalence relation in \(G\) defined by
\(G\times H\to G\times G\), \((g,h)\mapsto(g,gh)\) (cf. \SGARefLink{sga3:iv:subsection:3:2}{3.2}). The functor
\(Q\) defined by \(Q(S)=G(S)/H(S)\) is a quotient in
\(\widehat{\mathscr C}\) (by \SGARefLink{sga3:iv:proposition:4:4:9}{4.4.9} applied to the coarsest topology; cf.
\SGARefLink{sga3:iv:scholium:4:4:2}{4.4.2}), but it is not in general the quotient one wants. For example, when
\(\mathscr C=\mathbf{Sch}\), there is an exact sequence of affine group
schemes
\[
  1\longrightarrow\mu_2\longrightarrow\mathbf G_m
  \xrightarrow{p}\mathbf G_m\longrightarrow1
\]
that identifies \(\mathbf G_m\) with the quotient \(\mathbf G_m/\mu_2\).
Moreover, because \(p\) is finite and locally free, \(\mathbf G_m\) is the
sheaf quotient of \(\mathbf G_m\) by \(\mu_2\) in the larger category of
sheaves for the fppf topology; cf. \SGARefLink{sga3:iv:theorem:4:6:6:item:ii:01}{4.6.6(ii)} and \SGARefLink{sga3:iv:corollary:6:3:2}{6.3.2}. By contrast, the
quotient \(Q\) in \(\widehat{\mathscr C}\) is not isomorphic to
\(\mathbf G_m\), since, for example, \(Q(\mathbf Z)=\{1\}\), whereas
\(\mathbf G_m(\mathbf Z)=\{\pm1\}\). Thus \(Q\) is not an fppf sheaf and,
\emph{a fortiori}, is not representable.}

In what follows, we shall simply say \emph{equivalence relation} to mean a
\(\widehat{\mathscr C}\)-equivalence relation; whenever necessary, we shall
specify that a \(\mathscr C\)-equivalence relation is intended.

\medskip
\noindent\SGARefTarget{sga3:iv:definition:3:1:3}\textbf{Definition 3.1.3.}
\textit{If \(X\) is an object of \(\mathscr C\) over \(S\), an
\emph{equivalence relation in \(X\) over \(S\)} is an equivalence relation
\(R\) in \(X\) for which the structure morphism \(X\to S\) is compatible
with \(R\).}

The canonical monomorphism \(R\to X\times X\) then factors through the
monomorphism
\[
  X\times_S X\longrightarrow X\times X
\]
and defines an equivalence relation in the object \(X\to S\) of
\(\mathscr C_{/S}\). Whenever the quotient \(X/R\) exists, it is equipped
with a canonical morphism to \(S\), and the corresponding object of
\(\mathscr C_{/S}\) is a quotient of \(X\in\operatorname{Ob}\mathscr C_{/S}\)
by the preceding equivalence relation. Conversely, if \(S\) is a squarable
object of \(\mathscr C\) and \(Y\to S\) is a quotient of \(X\) by this
equivalence relation in \(\mathscr C_{/S}\), then \(Y\) is a quotient of
\(X\) by \(R\) in \(\mathscr C\). In any event, we shall never have to
consider quotients in \(\mathscr C_{/S}\) that are not already quotients in
\(\mathscr C\).

\medskip
\noindent\SGARefTarget{sga3:iv:definition:3:1:4}\textbf{Definition 3.1.4.}
\textit{Let \(X\) (respectively \(X'\)) be an object of \(\mathscr C\)
equipped with an equivalence relation \(R\) (respectively \(R'\)). A morphism
\[
  u:X\longrightarrow X'
\]
is said to be \emph{compatible with \(R\) and \(R'\)} if the following
equivalent conditions hold:}
\begin{enumerate}
  \item[(i)]\SGARefTarget{sga3:iv:definition:3:1:4:item:i:01} \textit{for every \(S\in\operatorname{Ob}\mathscr C\), two
  points of \(X(S)\) congruent modulo \(R(S)\) are carried by \(u\) to two
  points of \(X'(S)\) congruent modulo \(R'(S)\);}
  \item[(ii)]\SGARefTarget{sga3:iv:definition:3:1:4:item:ii:01} \textit{there exists a necessarily unique morphism \(R\to R'\)
  making the diagram}
  \[
    \begin{array}{ccc}
      R & \longrightarrow & R'\\
      \big\downarrow && \big\downarrow\\
      X\times X & \xrightarrow{u\times u} & X'\times X'
    \end{array}
  \]
  \textit{commute.}
\end{enumerate}

By the universal property of \(X/R\), whenever \(X/R\) and \(X'/R'\)
exist there is a unique morphism \(v\) making the diagram
\[
  \begin{array}{ccc}
    X & \xrightarrow{p} & X/R\\
    {\scriptstyle u}\big\downarrow && \big\downarrow{\scriptstyle v}\\
    X' & \xrightarrow{p'} & X'/R'
  \end{array}
\]
commute.

\medskip
\noindent\SGARefTarget{sga3:iv:definition:3:1:5}\textbf{Definition 3.1.5.}
\textit{A subobject (that is, a representable subfunctor) \(Y\) of \(X\) is
said to be \emph{stable under the equivalence relation \(R\)} if the following
equivalent conditions hold:}
\begin{enumerate}
  \item[(i)]\SGARefTarget{sga3:iv:definition:3:1:5:item:i:01} \textit{for every \(S\in\operatorname{Ob}\mathscr C\), the
  subset \(Y(S)\) of \(X(S)\) is stable under \(R(S)\);}
  \item[(ii)]\SGARefTarget{sga3:iv:definition:3:1:5:item:ii:01} \textit{the two subobjects of \(R\) that are the inverse
  images of \(Y\) under \(p_1\) and \(p_2\) are identical.}
\end{enumerate}

An important special case of a stable subobject of \(X\) is the following:
the quotient \(X/R\) exists and \(Y\) is the inverse image in \(X\) of a
subobject of \(X/R\).

\medskip
\noindent\SGARefTarget{sga3:iv:definition:3:1:6}\textbf{Definition 3.1.6.}
\textit{Let \(R\) be an equivalence relation in \(X\), and let \(X'\to X\)
be a morphism. The equivalence relation \(R'\) in \(X'\) obtained from the
cartesian diagram}
\[
  \begin{array}{ccc}
    R' & \longrightarrow & R\\
    \big\downarrow && \big\downarrow\\
    X'\times X' & \longrightarrow & X\times X
  \end{array}
\]
\textit{is called the equivalence relation in \(X'\) that is the
\emph{inverse image} of the equivalence relation \(R\) in \(X\). In
particular, if \(X'\) is a subobject of \(X\), it is called the equivalence
relation \emph{induced} on \(X'\) by \(R\), and is denoted \(R_{X'}\).}

The morphism \(X'\to X\) is compatible with \(R'\) and \(R\); hence, when
the quotients exist, it induces a morphism \(X'/R'\to X/R\) (\SGARefLink{sga3:iv:definition:3:1:4}{3.1.4}). If
\(X'\) is a subobject of \(X\), we shall see later that in certain cases one
can prove that \(X'/R'\to X/R\) is a monomorphism and therefore identifies
\(X'/R'\) with a subobject of \(X/R\). In that situation, the inverse image
of this subobject in \(X\) is a subobject containing \(X'\) and stable under
\(R\): the \emph{saturation} of \(X'\) with respect to the equivalence
relation \(R\).

\medskip
\noindent\SGARefTarget{sga3:iv:proposition:3:1:7}\textbf{Proposition 3.1.7.}
\textit{If the subobject \(Y\) of \(X\) is stable under \(R\), then for
\(i=1,2\) the following two squares are cartesian:}
\[
  \begin{array}{ccc}
    R_Y & \longrightarrow & R\\
    {\scriptstyle p_i}\big\downarrow && \big\downarrow{\scriptstyle p_i}\\
    Y & \longrightarrow & X.
  \end{array}
\]

The proof is immediate.

\SGARefTarget{sga3:iv:subsection:3:2}\subsection*{3.2. Equivalence Relation Defined by a Group Acting Freely}
\addcontentsline{toc}{subsection}{\SGARefLink{sga3:iv:subsection:3:2}{3.2}. Equivalence Relation Defined by a Group Acting Freely}

\noindent\SGARefTarget{sga3:iv:definition:3:2:1}\textbf{Definition 3.2.1.}
\textit{Let \(X\) be an object of \(\mathscr C\), and let \(H\) be a
\(\mathscr C\)-group acting on \(X\), that is, equipped with a morphism of
\(\widehat{\mathscr C}\)-groups}
\[
  H\longrightarrow\operatorname{Aut}(X).
\]
\textit{The group \(H\) is said to act \emph{freely} on \(X\) if the
following equivalent conditions hold:}
\begin{enumerate}
  \item[(i)]\SGARefTarget{sga3:iv:definition:3:2:1:item:i:01} \textit{for every \(S\in\operatorname{Ob}\mathscr C\), the
  group \(H(S)\) acts freely on \(X(S)\);}
  \item[(ii)]\SGARefTarget{sga3:iv:definition:3:2:1:item:ii:01} \textit{the morphism of functors}
  \[
    H\times X\longrightarrow X\times X
  \]
  \textit{defined on sets by \((h,x)\mapsto(hx,x)\) is a monomorphism.}
\end{enumerate}

In \SGARefLink{sga3:iv:definition:3:2:1}{the preceding definition} the group \(H\) was assumed to act on \(X\) on
the left. There is of course an analogous notion when the group acts on the
right, that is, when one is given a group morphism from the opposite group
\(H^{\circ}\) to \(\operatorname{Aut}(X)\).

If \(H\) acts freely on \(X\), the image of \(H\times X\) under the
morphism in \SGARefLink{sga3:iv:definition:3:2:1:item:ii:01}{(ii)} is an equivalence relation in \(X\), called the
\emph{equivalence relation defined by the action of \(H\) on \(X\)}. When
the quotient of \(X\) by this equivalence relation exists, it is denoted
\(H\backslash X\) (or \(X/H\) when \(H\) acts on the right). It represents
the following covariant functor: for every object \(Z\) of \(\mathscr C\),
\[
  \operatorname{Hom}(H\backslash X,Z)
  =\{\text{morphisms from \(X\) to \(Z\) invariant under \(H\)}\},
\]
where a morphism \(f:X\to Z\) is said to be invariant under \(H\) if, for
every \(S\in\operatorname{Ob}\mathscr C\), the corresponding morphism
\(X(S)\to Z(S)\) is invariant under the group \(H(S)\).

\medskip
\noindent\SGARefTarget{sga3:iv:lemma:3:2:2}\textbf{Lemma 3.2.2.}
\textit{Under the assumptions of \SGARefLink{sga3:iv:definition:3:2:1}{3.2.1}, let \(Y\) be a subobject of \(X\).
The following conditions are equivalent:}
\begin{enumerate}
  \item[(i)]\SGARefTarget{sga3:iv:lemma:3:2:2:item:i:01} \textit{\(Y\) is stable under the equivalence relation defined
  by \(H\) (\SGARefLink{sga3:iv:definition:3:1:5}{3.1.5});}
  \item[(ii)]\SGARefTarget{sga3:iv:lemma:3:2:2:item:ii:01} \textit{for every \(S\in\operatorname{Ob}\mathscr C\), the
  subset \(Y(S)\) of \(X(S)\) is stable under \(H(S)\);}
  \item[(iii)]\SGARefTarget{sga3:iv:lemma:3:2:2:item:iii:01} \textit{there exists a necessarily unique morphism \(f\)
  making the diagram}
  \[
    \begin{array}{ccc}
      H\times Y & \xrightarrow{f} & Y\\
      \big\downarrow && \big\downarrow\\
      H\times X & \longrightarrow & X
    \end{array}
  \]
  \textit{commute.}
\end{enumerate}
\textit{Under these conditions, \(f\) defines a morphism of
\(\widehat{\mathscr C}\)-groups}
\[
  H\longrightarrow\operatorname{Aut}(Y),
\]
\textit{and the equivalence relation in \(Y\) defined by this action of
\(H\) is precisely the equivalence relation induced on \(Y\) by the
equivalence relation in \(X\) defined by the action of \(H\).}

The proof is immediate. There is of course an analogous statement for a
right action. The action of \(H\) on \(Y\) defined above will be called the
action induced on \(Y\) by the given action of \(H\) on \(X\).

Consider now two \(\mathscr C\)-groups \(H\) and \(G\), together with a
group morphism
\[
  u:H\longrightarrow G.
\]
Then \(H\) acts on \(G\) by translations (set-theoretically,
\(hg=u(h)g\)), and it acts freely if and only if \(u\) is a monomorphism.
When it exists, the quotient of \(G\) by this action of \(H\) is denoted
\(H\backslash G\). Similarly, one defines a right action of \(H\) on \(G\)
and a quotient \(G/H\). These quotients are functorial with respect to the
groups involved. More precisely, one has \SGARefLink{sga3:iv:lemma:3:2:3}{the following lemma}, stated for
right quotients.

\medskip
\noindent\SGARefTarget{sga3:iv:lemma:3:2:3}\textbf{Lemma 3.2.3.}
\textit{Let \(u:H\to G\) and \(u':H'\to G'\) be monomorphisms of
\(\mathscr C\)-groups, and suppose that a morphism of \(\mathscr C\)-groups
\[
  f:G\longrightarrow G'
\]
is given. The following conditions are equivalent:}
\begin{enumerate}
  \item[(i)]\SGARefTarget{sga3:iv:lemma:3:2:3:item:i:01} \textit{\(f\) is compatible with the equivalence relations in
  \(G\) and \(G'\) defined by \(H\) and \(H'\);}
  \item[(ii)]\SGARefTarget{sga3:iv:lemma:3:2:3:item:ii:01} \textit{for every \(S\in\operatorname{Ob}\mathscr C\), one
  has \(fu(H(S))\subseteq u'(H'(S))\);}
  \item[(iii)]\SGARefTarget{sga3:iv:lemma:3:2:3:item:iii:01} \textit{there exists a necessarily unique multiplicative
  morphism \(g:H\to H'\) such that the diagram}
  \[
    \begin{array}{ccc}
      H & \xrightarrow{g} & H'\\
      {\scriptstyle u}\big\downarrow && \big\downarrow{\scriptstyle u'}\\
      G & \xrightarrow{f} & G'
    \end{array}
  \]
  \textit{commutes.}
\end{enumerate}
\textit{Under these conditions, if the quotients \(G/H\) and \(G'/H'\)
exist, there is a unique morphism \(\overline f\) making the diagram}
\[
  \begin{array}{ccc}
    G & \xrightarrow{f} & G'\\
    {\scriptstyle p}\big\downarrow && \big\downarrow{\scriptstyle p'}\\
    G/H & \xrightarrow{\overline f} & G'/H'
  \end{array}
\]
\textit{commute.}

The first assertion is proved by reduction to the set-theoretic case. The
second follows immediately from \SGARefLink{sga3:iv:lemma:3:2:3:item:i:01}{(i)}.

The notions introduced above for general equivalence relations could all be
translated into the present setting. We record only \SGARefLink{sga3:iv:lemma:3:2:4}{the following lemma},
whose proof is immediate by reduction to the set-theoretic case.

\medskip
\noindent\SGARefTarget{sga3:iv:lemma:3:2:4}\textbf{Lemma 3.2.4.}
\textit{Let \(u:H\to G\) be a monomorphism of \(\mathscr C\)-groups, and
let \(G'\) be a \(\mathscr C\)-subgroup of \(G\). The subobject \(G'\) of
\(G\) is stable under the equivalence relation defined by \(H\) if and only
if \(u\) factors through the canonical monomorphism \(G'\to G\). Under \SGARefLink{sga3:iv:lemma:3:2:4}{this
condition}, the action of \(H\) on \(G'\) induced by the action on \(G\)
defined by \(u\) is precisely the action arising from the monomorphism
\(H\to G'\) through which \(u\) factors.}

\SGARefTarget{sga3:iv:subsection:3:3}\subsection*{3.3. Universally Effective Equivalence Relations}
\addcontentsline{toc}{subsection}{\SGARefLink{sga3:iv:subsection:3:3}{3.3}. Universally Effective Equivalence Relations}

\noindent\SGARefTarget{sga3:iv:definition:3:3:1}\textbf{Definition 3.3.1.}
\textit{Let \(f:X\to Y\) be a morphism. The \emph{equivalence relation
defined by \(f\) in \(X\)}, denoted \(\mathscr R(f)\), is the
\(\widehat{\mathscr C}\)-equivalence relation in \(X\) that is the image of
the canonical monomorphism}
\[
  X\times_Y X\longrightarrow X\times X.
\]

\medskip
\noindent\SGARefTarget{sga3:iv:definition:3:3:2}\textbf{Definition 3.3.2.}
\textit{An equivalence relation \(R\) in \(X\) is said to be
\emph{effective} if:}
\begin{enumerate}
  \item[(i)]\SGARefTarget{sga3:iv:definition:3:3:2:item:i:01} \textit{\(R\) is representable, that is, it is a
  \(\mathscr C\)-equivalence relation;}
  \item[(ii)]\SGARefTarget{sga3:iv:definition:3:3:2:item:ii:01} \textit{the quotient \(Y=X/R\) exists in
  \(\mathscr C\);}
  \item[(iii)]\SGARefTarget{sga3:iv:definition:3:3:2:item:iii:01} \textit{the diagram}
  \[
    R\overset{p_1}{\underset{p_2}{\rightrightarrows}}
    X\xrightarrow{p}Y
  \]
  \textit{makes \(R\) the fiber square of \(X\) over \(Y\), that is,
  \(R\) is the equivalence relation defined by \(p\).}
\end{enumerate}

\medskip
\noindent\SGARefTarget{sga3:iv:scholium:3:3:2:1}\textbf{Scholium 3.3.2.1.}\footnote{\label{fnref:IV-14}Editor's note: the numbering
\SGARefLink{sga3:iv:scholium:3:3:2:1}{3.3.2.1} (respectively \SGARefLink{sga3:iv:scholium:3:3:3:1}{3.3.3.1}) has been added for later reference; Remark
\SGARefLink{sga3:iv:remark:3:3:3:2}{3.3.3.2} has also been added.}
\textit{If \(R\) is an effective equivalence relation in \(X\), then \(p\)
is an effective epimorphism (\SGARefLink{sga3:iv:definition:1:3}{1.3}). If \(f:X\to Y\) is an effective
epimorphism, then \(\mathscr R(f)\) is an effective equivalence relation in
\(X\), one quotient of which is \(Y\). Thus there is a one-to-one
``Galois'' correspondence between effective equivalence relations in \(X\)
and effective quotients of \(X\), that is, equivalence classes of effective
epimorphisms with source \(X\).}

\medskip
\noindent\SGARefTarget{sga3:iv:definition:3:3:3}\textbf{Definition 3.3.3.}
\textit{An equivalence relation \(R\) in \(X\) is said to be
\emph{universally effective} if the quotient \(Y=X/R\) exists and, for
every \(Y'\to Y\), the fiber products \(X'=X\times_Y Y'\) and
\(R'=R\times_Y Y'\) exist and \(R'\) is the fiber square of \(X'\) over
\(Y'\). Equivalently, \(R\) is effective and \(p:X\to X/R\) is a universal
effective epimorphism.}

\medskip
\noindent\SGARefTarget{sga3:iv:scholium:3:3:3:1}\textbf{Scholium 3.3.3.1.}\textsuperscript{\normalfont\ref{fnref:IV-14}}
\textit{As above, there is a one-to-one correspondence between universally
effective equivalence relations in \(X\) and universal effective quotients
of \(X\).}

\medskip
\noindent\SGARefTarget{sga3:iv:remark:3:3:3:2}\textbf{Remark 3.3.3.2.}\textsuperscript{\normalfont\ref{fnref:IV-14}}
Suppose that \(\mathscr C\) is the category of \(S\)-schemes, and write
\(\mathbf A^1_S\) for affine one-space over \(S\). Let
\(R\subset X\times_S X\) be a universally effective equivalence relation,
and let \(p:X\to Y\) be its quotient. Then, for every open subset \(U\) of
\(Y\),
\[
  \mathcal O(U)=\operatorname{Hom}_S(U,\mathbf A^1_S)
\]
is the set of elements \(\varphi\) of
\[
  \mathcal O(p^{-1}(U))
  =\operatorname{Hom}_S(p^{-1}(U),\mathbf A^1_S)
\]
such that \(\varphi\circ\operatorname{pr}_1
=\varphi\circ\operatorname{pr}_2\). In particular, if \(R\) is defined by
the action of a group \(H\) acting freely on the right on \(X\) (cf. \SGARefLink{sga3:iv:definition:3:2:1}{3.2.1}),
then \(\mathcal O(U)\) is the set of all
\(\varphi\in\mathcal O(p^{-1}(U))\) such that
\[
  \varphi(xh)=\varphi(x)
\]
for every \(S'\to S\), \(x\in X(S')\), and \(h\in H(S')\).

\medskip
\noindent\SGARefTarget{sga3:iv:proposition:3:3:4}\textbf{Proposition 3.3.4.}
\textit{Let \(R\) be a universally effective equivalence relation in
\(X\). Let \(f:X\to Z\) be a morphism compatible with \(R\), hence
factoring through a morphism \(g:X/R\to Z\). The following conditions are
equivalent:}
\begin{enumerate}
  \item[(i)]\SGARefTarget{sga3:iv:proposition:3:3:4:item:i:01} \textit{\(g\) is a monomorphism;}
  \item[(ii)]\SGARefTarget{sga3:iv:proposition:3:3:4:item:ii:01} \textit{\(R\) is the equivalence relation defined by \(f\).}
\end{enumerate}

Indeed, \SGARefLink{sga3:iv:proposition:3:3:4:item:i:01}{(i)} plainly implies \SGARefLink{sga3:iv:proposition:3:3:4:item:ii:01}{(ii)}, while the converse is precisely \SGARefLink{sga3:iv:corollary:2:5}{2.5}.

\medskip
\noindent\SGARefTarget{sga3:iv:definition:3:3:5}\textbf{Definition 3.3.5.}
\textit{Let \(H\) be a \(\mathscr C\)-group acting freely on \(X\). One
says that \(H\) acts \emph{effectively}, or that the action of \(H\) on
\(X\) is \emph{effective} (respectively \emph{universally effective}), if
the equivalence relation in \(X\) defined by the action of \(H\) is
effective (respectively universally effective).}

\SGARefTarget{sga3:iv:subsection:3:4}\subsection*{3.4. \((M)\)-Effectivity}
\addcontentsline{toc}{subsection}{\SGARefLink{sga3:iv:subsection:3:4}{3.4}. (M)-Effectivity}

In practice, universal effective epimorphisms are most often difficult to
characterize. Nevertheless, one frequently has at hand a number of
morphisms of this kind--in scheme theory, for example, faithfully flat
quasi-compact morphisms. This leads to the developments below.

\medskip
\noindent\SGARefTarget{sga3:iv:para:3:4:1}\textbf{3.4.1.}
We first state several conditions on a family \((M)\) of morphisms of
\(\mathscr C\):
\begin{enumerate}
  \item[(a)]\SGARefTarget{sga3:iv:para:3:4:1:item:a:01} \((M)\) is stable under base change: every
  \(u:T\to S\) belonging to \((M)\) is squarable (cf. \SGARefLink{sga3:iv:definition:1:4:0}{1.4.0}), and for every
  \(S'\to S\), the base change
  \[
    u':T\times_S S'\longrightarrow S'
  \]
  belongs to \((M)\);
  \item[(b)]\SGARefTarget{sga3:iv:para:3:4:1:item:b:01} the composite of two elements of \((M)\) belongs to \((M)\);
  \item[(c)]\SGARefTarget{sga3:iv:para:3:4:1:item:c:01} every isomorphism belongs to \((M)\);
  \item[(d)]\SGARefTarget{sga3:iv:para:3:4:1:item:d:01} every element of \((M)\) is an effective epimorphism.
\end{enumerate}

Conditions \SGARefLink{sga3:iv:para:3:4:1:item:a:01}{(a)} and \SGARefLink{sga3:iv:para:3:4:1:item:b:01}{(b)} imply:
\begin{enumerate}
  \item[(a$'$)]\SGARefTarget{sga3:iv:para:3:4:1:item:a:prime:01} the cartesian product of two elements of \((M)\) belongs
  to \((M)\). More precisely, let \(u:X\to Y\) and \(u':X'\to Y'\) be two
  \(S\)-morphisms belonging to \((M)\). If \(Y\times_S Y'\) exists, then
  \(X\times_S X'\) exists and \(u\times_Su'\) belongs to \((M)\).
\end{enumerate}

This follows from the diagram
\[
\begin{array}{ccccc}
  &&Y'&\xleftarrow{u'}&X'\\
  &&\big\uparrow&&\big\uparrow\\[-2pt]
  X&\longleftarrow&X\times_S Y'&\longleftarrow&X\times_S X'\\
  {\scriptstyle u}\big\downarrow&&\big\downarrow&&
    \underset{\scriptstyle u\times_Su'}{\swarrow}\\[-2pt]
  Y&\longleftarrow&Y\times_S Y'&&
\end{array}
\]
Likewise, \SGARefLink{sga3:iv:para:3:4:1:item:a:01}{(a)} and \SGARefLink{sga3:iv:para:3:4:1:item:d:01}{(d)} imply:
\begin{enumerate}
  \item[(d$'$)]\SGARefTarget{sga3:iv:para:3:4:1:item:d:prime:01} every element of \((M)\) is a universal effective
  epimorphism.
\end{enumerate}

\medskip
\noindent\SGARefTarget{sga3:iv:para:3:4:2}\textbf{3.4.2.}
The family \((M_0)\) of universal effective epimorphisms satisfies
conditions \SGARefLink{sga3:iv:para:3:4:1:item:a:01}{(a)}--\SGARefLink{sga3:iv:para:3:4:1:item:d:01}{(d)} of \SGARefLink{sga3:iv:para:3:4:1}{3.4.1}. Indeed, \SGARefLink{sga3:iv:para:3:4:1:item:a:01}{(a)}, \SGARefLink{sga3:iv:para:3:4:1:item:c:01}{(c)}, and \SGARefLink{sga3:iv:para:3:4:1:item:d:01}{(d)} hold by definition,
and \SGARefLink{sga3:iv:para:3:4:1:item:b:01}{(b)} follows from \SGARefLink{sga3:iv:proposition:1:8}{1.8}. Henceforth we assume given a family \((M)\) of
morphisms of \(\mathscr C\) satisfying these conditions; our results apply
in particular to the family \((M_0)\).

\medskip
\noindent\SGARefTarget{sga3:iv:definition:3:4:3}\textbf{Definition 3.4.3.}
An equivalence relation \(R\) in \(X\) is said to be \emph{of type \((M)\)}
if it is representable and \(p_1\in(M)\); conditions \SGARefLink{sga3:iv:para:3:4:1:item:b:01}{(b)} and \SGARefLink{sga3:iv:para:3:4:1:item:c:01}{(c)} then imply
that \(p_2\in(M)\).

The relation \(R\) is said to be \emph{\((M)\)-effective} if it is effective
and the canonical morphism \(X\to X/R\) belongs to \((M)\).

A quotient \(Y\) of \(X\) is said to be \emph{\((M)\)-effective} if the
canonical morphism \(X\to Y\) belongs to \((M)\).

The following consequences are immediate from \SGARefLink{sga3:iv:definition:3:4:3}{the definition}.\footnote{Editor's
note: the numbering \SGARefLink{sga3:iv:proposition:3:4:3:1}{3.4.3.1} has been added for later reference, and the proof
of part \SGARefLink{sga3:iv:proposition:3:4:3:1:item:i:01}{(i)} has been expanded.}

\medskip
\noindent\SGARefTarget{sga3:iv:proposition:3:4:3:1}\textbf{Proposition 3.4.3.1.}
\begin{enumerate}
  \item[(i)]\SGARefTarget{sga3:iv:proposition:3:4:3:1:item:i:01} \textit{An \((M)\)-effective equivalence relation is of type
  \((M)\) and is universally effective.}
  \item[(ii)]\SGARefTarget{sga3:iv:proposition:3:4:3:1:item:ii:01} \textit{An \((M)\)-effective quotient is universally effective
  (cf. \SGARefLink{sga3:iv:definition:3:3:3}{3.3.3}).}
  \item[(iii)]\SGARefTarget{sga3:iv:proposition:3:4:3:1:item:iii:01} \textit{The maps \(R\mapsto X/R\) and
  \(p\mapsto\mathscr R(p)\) establish a one-to-one correspondence between
  \((M)\)-effective equivalence relations in \(X\) and
  \((M)\)-effective quotients of \(X\).}
  \item[(iv)]\SGARefTarget{sga3:iv:proposition:3:4:3:1:item:iv:01} \textit{To be \((M_0)\)-effective is equivalent to being
  universally effective.}
\end{enumerate}

We prove \SGARefLink{sga3:iv:proposition:3:4:3:1:item:i:01}{(i)}. Since \(R\) is \((M)\)-effective, there is a cartesian square
\[
  \begin{array}{ccc}
    R & \xrightarrow{p_2} & X\\
    {\scriptstyle p_1}\big\downarrow && \big\downarrow{\scriptstyle p}\\
    X & \xrightarrow{p} & X/R,
  \end{array}
\]
and \(p\in(M)\). By \SGARefLink{sga3:iv:para:3:4:1:item:a:01}{3.4.1(a)}, \(p_1\) and \(p_2\) belong to \((M)\), so
\(R\) is of type \((M)\).

Set \(Y=X/R\), and let \(Y'\to Y\) be any morphism. By \SGARefLink{sga3:iv:para:3:4:1:item:a:01}{3.4.1(a)}, the fiber
products \(X'=X\times_Y Y'\) and \(R'=R\times_Y Y'\) exist, and the
morphisms \(X'\to Y'\) and \(p_i':R'\to X'\) belong to \((M)\). Finally,
since \(R=X\times_Y X\), associativity of the fiber product gives
\[
  R'=X\times_Y X\times_Y Y'=X'\times_{Y'}X'.
\]
This shows that \(R'\) is \((M)\)-effective; in particular, \(R\) is
universally effective. This proves \SGARefLink{sga3:iv:proposition:3:4:3:1:item:i:01}{(i)}, as well as \SGARefLink{sga3:iv:proposition:3:4:3:1:item:iv:01}{(iv)}. Parts \SGARefLink{sga3:iv:proposition:3:4:3:1:item:ii:01}{(ii)} and \SGARefLink{sga3:iv:proposition:3:4:3:1:item:iii:01}{(iii)}
follow, taking Definition \SGARefLink{sga3:iv:definition:3:3:2}{3.3.2} into account.

\medskip
\noindent\SGARefTarget{sga3:iv:para:3:4:4}\textbf{3.4.4.}
Let \(H\) be an \(S\)-group whose structure morphism belongs to \((M)\). If
\(H\) acts freely on the \(S\)-object \(X\), then it defines an equivalence
relation of type \((M)\) in \(X\). Indeed, by \SGARefLink{sga3:iv:para:3:4:1:item:a:01}{(a)}, the fiber product
\(H\times_S X\) exists and
\(\operatorname{pr}_2:H\times_S X\to X\) belongs to \((M)\). The action of
\(H\) is said to be \emph{\((M)\)-effective} if the equivalence relation in
\(X\) defined by this action is \((M)\)-effective.

\medskip
\noindent\SGARefTarget{sga3:iv:proposition:3:4:5}\textbf{Proposition 3.4.5 (\((M)\)-effectivity and base change).}
\textit{Let \(R\) be an \((M)\)-effective equivalence relation in \(X\)
over \(S\), and set \(Y=X/R\). Let \(S'\to S\) be a base change for which
\(Y'=Y\times_S S'\) exists. Then \(X'=X\times_S S'\) exists,
\(R'=R\times_S S'\) is an \((M)\)-effective equivalence relation in \(X'\)
over \(S'\), and}
\[
  X'/R'\simeq (X/R)'.
\]

Indeed, the canonical morphisms \(X\to Y\) and \(R\to Y\) belong to
\((M)\); hence, by \SGARefLink{sga3:iv:para:3:4:1:item:a:prime:01}{(a$'$)}, \(X'\) and \(R'\) are representable. By
associativity of the fiber product, \(R'\) is the equivalence relation in
\(X'\) defined by the canonical morphism \(X'\to Y'\), which belongs to
\((M)\). This proves \SGARefLink{sga3:iv:proposition:3:4:5}{the assertion}.

\medskip
\noindent\SGARefTarget{sga3:iv:proposition:3:4:6}\textbf{Proposition 3.4.6 (\((M)\)-effectivity and cartesian
products).}
\textit{Let \(R\) (respectively \(R'\)) be an \((M)\)-effective equivalence
relation in \(X\) (respectively \(X'\)) over \(S\). If
\((X/R)\times_S(X'/R')\) exists, then \(X\times_SX'\) exists,
\(R\times_SR'\) is an \((M)\)-effective equivalence relation in
\(X\times_SX'\) over \(S\), and}
\[
  (X\times_SX')/(R\times_SR')
  \simeq (X/R)\times_S(X'/R').
\]

Set \(Y=X/R\) and \(Y'=X'/R'\). By \SGARefLink{sga3:iv:para:3:4:1:item:a:prime:01}{(a$'$)}, the fiber product
\(X\times_SX'\) exists and the canonical morphism
\[
  q:X\times_SX'\longrightarrow Y\times_SY'
\]
belongs to \((M)\). The formula
\[
  (X\times_SX')\times_{(Y\times_SY')}(X\times_SX')
  \simeq (X\times_YX)\times_S(X'\times_{Y'}X')
\]
(all unlabelled fiber products being taken over \(S\)) shows that
\(R\times_SR'\) is the equivalence relation in \(X\times_SX'\) defined by
\(q\), which completes the proof.

\medskip
\noindent\SGARefTarget{sga3:iv:para:3:4:7}\textbf{3.4.7.}\footnote{Editor's note: \SGARefLink{sga3:iv:para:3:4:7}{this paragraph} has been
added.}
Suppose that \(\mathscr C\) has a final object \(e\), and let
\(f:G\to G'\) be a morphism of \(\mathscr C\)-groups belonging to \((M)\).
By \SGARefLink{sga3:iv:para:3:4:1:item:a:01}{3.4.1(a)}, the kernel \(\operatorname{Ker}(f)\) is represented by
\(e\times_{G'}G\), and the morphism \(\operatorname{Ker}(f)\to e\) belongs
to \((M)\).

On the other hand, the equivalence relation defined by \(f\) is the same as
that defined by the action of \(\operatorname{Ker}(f)\), say on the right,
on \(G\): it is the image of the morphism
\[
  G\times\operatorname{Ker}(f)\longrightarrow G\times G
\]
defined on sets by \((g,h)\mapsto(g,gh)\). Consequently, Scholium \SGARefLink{sga3:iv:scholium:3:3:2:1}{3.3.2.1}
gives \SGARefLink{sga3:iv:corollary:3:4:7:1}{the following corollary}.

\medskip
\noindent\SGARefTarget{sga3:iv:corollary:3:4:7:1}\textbf{Corollary 3.4.7.1.}
\textit{Suppose that \(\mathscr C\) has a final object \(e\), and let
\(f:G\to G'\) be a morphism of \(\mathscr C\)-groups belonging to \((M)\).
Then the action of \(\operatorname{Ker}(f)\) on \(G\) is
\((M)\)-effective, and \(G'\) is the quotient
\(G/\operatorname{Ker}(f)\).}

\SGARefTarget{sga3:iv:subsection:3:5}\subsection*{3.5. Construction of Quotients by Descent}
\addcontentsline{toc}{subsection}{\SGARefLink{sga3:iv:subsection:3:5}{3.5}. Construction of Quotients by Descent}

It frequently happens that one cannot construct a quotient directly, but
can do so after a suitable base change. This subsection gives a useful
criterion for that situation.

In \SGARefLink{sga3:iv:definition:2:1}{2.1} we defined a descent datum on an object \(X'\) over \(S'\), relative
to a morphism \(S'\to S\).

\medskip
\noindent\SGARefTarget{sga3:iv:definition:3:5:1}\textbf{Definition 3.5.1.}
\textit{A descent datum on \(X'\), relative to \(S'\to S\), is said to be
\emph{effective} if \(X'\), equipped with this descent datum, is isomorphic
to the inverse image over \(S'\) of an object \(X\) over \(S\), equipped
with its canonical descent datum.}

If \(S'\to S\) is a descent morphism (\SGARefLink{sga3:iv:definition:2:2}{2.2}), then the object \(X\) in \SGARefLink{sga3:iv:definition:3:5:1}{the
definition} is unique up to unique isomorphism. To say that \(S'\to S\) is
an effective descent morphism (\SGARefLink{sga3:iv:definition:2:2}{2.2}) means that it is a descent morphism and
that every descent datum relative to it is effective.

Now consider an equivalence relation \(R\) in an object \(X\) over \(S\).
Let \(X'\), \(X''\), and \(X'''\) be the inverse images of \(X\) over
\[
  S',\qquad S''=S'\times_SS',\qquad
  S'''=S'\times_SS'\times_SS',
\]
respectively, and let \(R'\), \(R''\), and \(R'''\) be the equivalence
relations obtained from \(R\) by inverse image. Suppose that \(R'\) in
\(X'\) is \((M)\)-effective, and consider the quotient
\(Y'=X'/R'\), an object over \(S'\). By \SGARefLink{sga3:iv:proposition:3:4:5}{3.4.5}, its two inverse images over
\(S''\) are isomorphic to \(X''/R''\). Thus the \(S'\)-object \(Y'\) has a
canonical gluing datum. Applying in the same way the uniqueness of
\(X'''/R'''\), one sees that this is in fact a descent datum. (In this
argument it has been tacitly assumed that all the displayed fiber products
exist. This holds, in particular, if \(S'\to S\) is squarable, for example
if it is a descent morphism.)

\medskip
\noindent\SGARefTarget{sga3:iv:proposition:3:5:2}\textbf{Proposition 3.5.2.}
\textit{Let \(R\) be an equivalence relation in the object \(X\) over
\(S\), and let \(S'\to S\) be a universal effective epimorphism. Assume
that every \(S\)-morphism whose inverse image over \(S'\) belongs to
\((M)\) itself belongs to \((M)\). The following conditions are equivalent:}
\begin{enumerate}
  \item[(i)]\SGARefTarget{sga3:iv:proposition:3:5:2:item:i:01} \textit{\(R\) is \((M)\)-effective in \(X\);}
  \item[(ii)]\SGARefTarget{sga3:iv:proposition:3:5:2:item:ii:01} \textit{\(R'\) is \((M)\)-effective in \(X'\), and the
  canonical descent datum on \(X'/R'\) is effective.}
\end{enumerate}
\textit{When these conditions hold, the object descended from \(X'/R'\) is
canonically isomorphic to \(X/R\).}

The implication \SGARefLink{sga3:iv:proposition:3:5:2:item:i:01}{(i)}\(\Rightarrow\)\SGARefLink{sga3:iv:proposition:3:5:2:item:ii:01}{(ii)} is simply \SGARefLink{sga3:iv:para:3:4:4}{3.4.4} translated into the
language of descent. Once the converse is proved, the last assertion of \SGARefLink{sga3:iv:proposition:3:5:2}{the
proposition} follows from the fact that a universal effective epimorphism is
a descent morphism (\SGARefLink{sga3:iv:proposition:2:3}{2.3}), and hence the descended object is unique up to
unique isomorphism.

We prove \SGARefLink{sga3:iv:proposition:3:5:2:item:ii:01}{(ii)}\(\Rightarrow\)\SGARefLink{sga3:iv:proposition:3:5:2:item:i:01}{(i)}. Let \(Y'=X'/R'\), and let \(Y\) be the
descended object. The canonical morphism
\(p':X'\to X'/R'=Y'\) is compatible by construction with the descent data:
its two inverse images over \(S''\) coincide with the canonical morphism
\(X''\to X''/R''\). It is therefore the inverse image over \(S'\) of an
\(S\)-morphism \(p:X\to Y\). Since \(p'\) belongs to \((M)\), the assumption
on \(S'\to S\) implies that \(p\) also belongs to \((M)\). Moreover,
because \(p'\) is compatible with the equivalence relation \(R'\), the
morphism \(p\) is compatible with \(R\), again because a universal
effective epimorphism is a descent morphism. Hence there is a morphism
\[
  R\longrightarrow X\times_YX.
\]
To prove that \(R\) is \((M)\)-effective and that \(Y\) is isomorphic to
\(X/R\), it suffices to prove that this morphism is an isomorphism. After
base change from \(S\) to \(S'\), it becomes an isomorphism because \(R'\)
is effective; it is therefore an isomorphism for the same reason as in \SGARefLink{sga3:iv:proposition:2:4}{2.4}.

Notice that the hypothesis in \SGARefLink{sga3:iv:proposition:3:5:2}{the proposition} is satisfied if \((M)\) is
the family \((M_0)\) of universal effective epimorphisms and if
\(\mathscr C\) has fiber products (\SGARefLink{sga3:iv:corollary:1:10}{1.10}). We obtain the following.

\medskip
\noindent\SGARefTarget{sga3:iv:corollary:3:5:3}\textbf{Corollary 3.5.3.}
\textit{Suppose that \(\mathscr C\) has fiber products (fiber products over
\(S\) would suffice). Let \(R\) be an equivalence relation in \(X\) over
\(S\), and let \(S'\to S\) be a universal effective epimorphism. The
following conditions are equivalent:}
\begin{enumerate}
  \item[(i)]\SGARefTarget{sga3:iv:corollary:3:5:3:item:i:01} \textit{\(R\) is universally effective in \(X\);}
  \item[(ii)]\SGARefTarget{sga3:iv:corollary:3:5:3:item:ii:01} \textit{\(R'\) is universally effective in \(X'\), and the
  canonical descent datum on \(X'/R'\) is effective.}
\end{enumerate}
\textit{When these conditions hold, the object descended from \(X'/R'\) is
canonically isomorphic to \(X/R\).}

